\chapter{Embauche et rémunération}\label{ch:fm:hiring-and-compensation}

Un trader rémunéré à hauteur de quinze pour cent d'une année à 20~millions de dollars et un trader qui ne touche rien après une perte de 2~millions de dollars sur le desk voisin peuvent avoir
exactement le même talent. Pour un trader dont le risque donne une volatilité annuelle du P\&L de 10~millions de dollars et dont le vrai ratio de Sharpe est égal à la
moyenne du desk, soit 0.8, la première année survient environ une fois sur neuf et la seconde une fois sur six. Une année de P\&L mesure bien plus la chance que la compétence~: sur
le desk du chapitre, 87~\% des écarts de \omterm{def:fm:hiring-and-compensation:pool}{rémunération à la formule} entre traders sur une même année relèvent de la chance. Une règle de rémunération décide laquelle des deux l'entreprise est en train
d'acheter, et combien de temps il faut pour le savoir.

\section{Le processus de recrutement et l'entretien}

Une société de trading recrute à partir d'un entonnoir~: candidatures, présélections, tests, entretiens, offres, acceptations. Deux faits le façonnent. Le taux de base est faible, si bien qu'un
test à taux modeste de faux positifs remplit encore les étapes suivantes de candidats qui ne seront pas recrutés~; chaque décision d'étape est une mesure
bruitée, donc l'ordre des étapes compte~: des tests peu coûteux et valides d'abord, les plus chers en dernier.

\begin{definition}[Entretien structuré]\label{def:fm:hiring-and-compensation:structured}
Un \emph{entretien structuré}\index{entretien structuré} pose à tous les candidats d'un même poste les mêmes questions, choisies à partir d'une analyse du poste,
et note les réponses sur des échelles ancrées convenues à l'avance, plusieurs examinateurs notant de façon indépendante avant de se concerter.
\end{definition}

La recherche sur la sélection est étonnamment tranchée sur un point. La synthèse de Schmidt et Hunter sur 85 ans d'études (1998) et sa révision par
Sackett, Zhang, Berry et Lievens (2022), qui ont constaté que les estimations antérieures avaient été surestimées de 0.10 à 0.20 par suite de surcorrections, s'accordent
à dire que les \omterm{def:fm:hiring-and-compensation:structured}{entretiens structurés} prédisent mieux la performance professionnelle que les entretiens non structurés~; dans la révision, ils sont la procédure de sélection la mieux classée.
Les questions d'entretien de cette série sont écrites dans ce but~: la même question, une bonne réponse connue, et ce que l'examinateur
recherche. Les livres~18 à~20 sont des recueils d'entraînement de nouvelles questions dans cet esprit, chacun avec une réponse complète.

\begin{method}[Concevoir un parcours d'entretiens]\label{met:fm:hiring-and-compensation:loop}
\begin{enumerate}
\item Rédiger l'analyse du poste~: les trois à cinq aptitudes qui distinguent les bons éléments des autres.
\item Pour chaque aptitude, choisir un test ou des questions structurées, avec des notes ancrées.
\item Placer en premier les tests peu coûteux et valides (une série de problèmes chronométrée, un test de programmation) et en dernier les plus coûteux (les parcours sur site).
\item Noter de façon indépendante avant de discuter~; consigner les notes et les résultats, et vérifier un an plus tard quelles notes ont prédit la performance.
\end{enumerate}
\end{method}

\section{Structures de rémunération~: fixe, bonus et différé}

La rémunération dans le trading comporte trois volets~: le salaire de base, fixé pour l'année~; un bonus, décidé après coup~; et, au-delà d'un certain niveau, une part différée qui s'acquiert
sur les années suivantes. L'équilibre entre ces éléments relève du choix de l'entreprise là où aucune règle ne le limite, et de celui d'un régulateur là où une règle le fait.

\begin{definition}[Enveloppe de bonus, rémunération à la formule, bonus discrétionnaire]\label{def:fm:hiring-and-compensation:pool}
Une \emph{enveloppe de bonus}\index{enveloppe de bonus} (bonus pool) est la rémunération variable totale qu'une entreprise ou un desk attribue pour une année, fixée par une règle ou par une décision à partir de ses résultats.
Une \emph{rémunération à la formule}\index{rémunération à la formule} (formulaic payout) verse à un trader ou à une équipe une part fixe de leur propre P\&L, généralement nette des coûts et des pertes
reportées (chapitre~3). Un \emph{bonus discrétionnaire}\index{bonus discrétionnaire} (discretionary bonus) est fixé selon le jugement de la direction
sur la contribution individuelle, au sein de l'enveloppe.
\end{definition}

\begin{definition}[Rémunération différée, calendrier d'acquisition]\label{def:fm:hiring-and-compensation:deferral}
La \emph{rémunération différée}\index{rémunération différée} (deferred compensation) est la rémunération variable attribuée pour une année mais versée les années suivantes, en numéraire, en actions ou en parts de
fonds. Son \emph{calendrier d'acquisition}\index{calendrier d'acquisition} (vesting schedule) fixe le moment où chaque part devient acquise au salarié~: en bloc à la fin d'une période (cliff) ou
par parts égales chaque année (pro rata).
\end{definition}

\begin{definition}[Malus, récupération (clawback)]\label{def:fm:hiring-and-compensation:malus}
Le \emph{malus}\index{malus} est la réduction ou l'annulation de la rémunération variable attribuée mais non encore acquise, en cas de faute, d'une erreur significative
ou de pertes apparues plus tard. La \emph{récupération}\index{récupération} (clawback) est le recouvrement de la rémunération variable déjà versée ou
acquise.
\end{definition}

\begin{definition}[Preneur de risques significatif]\label{def:fm:hiring-and-compensation:mrt}
Un \emph{preneur de risques significatif}\index{preneur de risques significatif} (material risk taker) est un salarié dont les activités, au titre des règles de rémunération applicables aux banques et aux entreprises d'investissement,
ont un effet significatif sur le profil de risque de l'établissement~; leur rémunération variable est soumise aux règles relatives au différé, aux instruments, au \omterm{def:fm:hiring-and-compensation:malus}{malus} et
à la \omterm{def:fm:hiring-and-compensation:malus}{récupération}.
\end{definition}

\begin{dated}{2026-09}{Les règles de rémunération}\label{dat:fm:hiring-and-compensation:rules}
\textbf{Union européenne} (Directive 2013/36/EU, article~94, telle que modifiée par la Directive (EU) 2019/878)~: pour les \omterm{def:fm:hiring-and-compensation:mrt}{preneurs de risques significatifs} des banques, la rémunération variable
ne peut dépasser 100~\% de la rémunération fixe, ou 200~\% avec l'approbation des actionnaires~; au moins la moitié est versée en actions ou en instruments équivalents~; au moins 40~\%
est différée sur quatre à cinq ans, avec une acquisition pas plus rapide que pro rata (60~\% pour les montants particulièrement élevés)~; jusqu'à 100~\% est soumise au \omterm{def:fm:hiring-and-compensation:malus}{malus} ou
à la \omterm{def:fm:hiring-and-compensation:malus}{récupération}~; une rémunération variable garantie n'est permise qu'à l'embauche, pour la première année. \textbf{Royaume-Uni}~: la PRA et la FCA ont supprimé la
limite de ratio (le «~plafond des bonus~») pour les banques, les building societies et les entreprises d'investissement désignées par la PRA, à compter du 31 octobre 2023 (PRA PS9/23, FCA PS23/15), en invoquant
son effet de hausse de la rémunération fixe. \textbf{États-Unis}~: la Rule 10D-1 de la SEC impose aux sociétés cotées de récupérer les rémunérations incitatives attribuées par erreur aux
dirigeants au cours des trois exercices qui précèdent un retraitement comptable. \textbf{Hauts revenus}~: l'Autorité bancaire européenne (EBA) a dénombré 2\,554
personnes gagnant plus de 1~million d'euros dans les banques et les entreprises d'investissement de l'UE en 2024, contre 2\,343 en 2023~; dans les entreprises d'investissement, le nombre a progressé de 30~\%,
de 221 à 288.
\end{dated}

\section{Dimensionner une enveloppe de bonus}

Une enveloppe se dimensionne par le haut, puis se répartit. Les règles courantes sont une part du revenu (le \omterm{def:fm:the-economics-of-a-trading-firm:comp}{ratio de rémunération} du chapitre~1, fixé à l'avance) ou une
part du résultat net, déduction faite des coûts et d'une \omterm{def:fm:limits-and-capital-allocation:raroc}{charge en capital} (chapitre~7). La seconde ne rémunère que la valeur créée~: un desk qui gagne moins que son
coût du capital n'a pas d'enveloppe, quel que soit son revenu. Le choix compte surtout les années où l'entreprise souhaite le moins payer~: une enveloppe fondée sur le revenu verse
des sommes une année où le desk a détruit de la valeur sur son capital.

\begin{example}[Deux enveloppes]\label{ex:fm:hiring-and-compensation:pools}
Un desk réalise un revenu de 100 avec des coûts de 60 et une \omterm{def:fm:limits-and-capital-allocation:raroc}{charge en capital} de 10. Une enveloppe de 20~\% du résultat après la charge vaut $0.2\times(100-60-10)=6$. Une
enveloppe de 6~\% du revenu vaut aussi 6 cette année. Si le revenu vaut 70 une année, la première enveloppe est nulle et la seconde vaut 4.2.
\end{example}

\section{Rémunération à la formule contre bonus discrétionnaire~: payer le talent, pas la chance}

Le P\&L d'un trader sur une année est la somme du talent et de la chance. Si le talent, pris sur l'ensemble des traders, a un écart-type $s$ (en P\&L par an) et la chance, c'est-à-dire le bruit de l'année, un
écart-type $n$, alors après $T$ années le P\&L moyen est corrélé au talent à $\sqrt{s^2/(s^2+n^2/T)}$.

\begin{proposition}[Combien de temps il faut à la rémunération pour révéler le talent]\label{prop:fm:hiring-and-compensation:years}
Soit le P\&L annuel du trader $i$ égal à $\mu_i+\varepsilon_{i,t}$, avec $\mu_i$ tiré avec variance $s^2$ et $\varepsilon_{i,t}$ indépendant de
variance $n^2$. La part de la variance transversale du P\&L d'une année due à la chance vaut $n^2/(s^2+n^2)$, et la corrélation de la moyenne sur $T$ années
du P\&L avec $\mu_i$ atteint $\rho$ après
\[
T=\frac{n^2}{s^2\,(1/\rho^2-1)}\ \text{ans}.
\]
\end{proposition}

\begin{proof}
La moyenne vaut $\mu_i+\bar\varepsilon_i$ avec $\mathrm{Var}(\bar\varepsilon_i)=n^2/T$~; sa covariance avec $\mu_i$ vaut $s^2$, donc sa corrélation
vaut $s/\sqrt{s^2+n^2/T}$. En égalant à $\rho$ et en résolvant, on obtient $T$.
\end{proof}

Sur le desk du chapitre (à titre d'illustration~: vingt traders, une volatilité annuelle du P\&L de 10~millions de dollars chacun, des ratios de Sharpe vrais dispersés avec un écart-type
de 0.4, soit $s=4$ et $n=10$), la chance représente 86~\% de la variance du P\&L d'une année, et il faut 11.1 ans de P\&L pour que celui-ci soit corrélé au talent à 0.8.
Une \omterm{def:fm:hiring-and-compensation:pool}{rémunération à la formule} de 15~\% du P\&L positif hérite des deux~: 87~\% de la variance de la rémunération d'une année relève de la chance, et sa corrélation avec le talent vaut 0.35 après
un an, 0.64 après cinq ans et 0.80 après environ douze ans (la \cref{fig:fm:hiring-and-compensation:years}).

\begin{omfigure}
\begin{tikzpicture}
\begin{axis}[width=10.5cm,height=5.2cm,xlabel={\small années de P\&L},ylabel={\small corrélation de la rémunération avec le talent},xmin=1,xmax=15,ymin=0,ymax=1,
  tick label style={font=\footnotesize},grid=major,grid style={gray!20},table/col sep=comma,
  legend cell align=left,legend style={font=\scriptsize,at={(0.5,-0.3)},anchor=north,/tikz/every even column/.append style={column sep=8pt}},legend columns=2]
\addplot[omDef,thick,mark=*,mark size=1.2pt] table[x=years,y=simulated]{figdata/desk/10-hiring-and-compensation/corr.csv};
\addplot[omThm,thick,dashed] table[x=years,y=linear]{figdata/desk/10-hiring-and-compensation/corr.csv};
\draw[gray,dotted] (axis cs:1,0.8) -- (axis cs:15,0.8);
\legend{rémunération à la formule cumulée (simulée),{P\&L moyen (\cref{prop:fm:hiring-and-compensation:years})}}
\end{axis}
\end{tikzpicture}

\omcaption{Le temps que met la rémunération à refléter le talent~: la corrélation, sur vingt traders, entre le ratio de Sharpe vrai et la rémunération cumulée, dans le cas d'une
\omterm{def:fm:hiring-and-compensation:pool}{rémunération à la formule} de 15~\% (400 desks simulés), et entre le talent et le P\&L moyen en forme close. Dispersion du talent de 0.4 en ratio de Sharpe, volatilité annuelle
10 millions de dollars. Données~: \texttt{fm\_pay.corr\_by\_years}.}\label{fig:fm:hiring-and-compensation:years}
\end{omfigure}

Une règle discrétionnaire ne peut faire mieux que si elle utilise une information que le P\&L ne contient pas~: la qualité des décisions du trader, sa prise de risque,
sa contribution au P\&L des autres. Le jugement appliqué au seul P\&L n'est au mieux qu'une estimation par rétrécissement du talent
(\cref{lst:fm:hiring-and-compensation:shrink}), et sur le desk du chapitre elle a une corrélation avec le talent de 0.37 après un an, contre 0.35 pour la formule,
et un peu moins après cinq ans (0.61 contre 0.64), parce que le rétrécissement est proportionnel et que l'enveloppe partagée suit la chance du desk. Une
rémunération ajustée du risque, qui fait payer à chaque trader sa \omterm{def:fm:limits-and-capital-allocation:raroc}{charge en capital}, présente une corrélation de 0.34 et de 0.63~: elle rémunère la valeur ajoutée, non la détection du talent.

\omcode{code/firm/bonuspool/firm_bonuspool.py}{44}{56}{Une allocation discrétionnaire qui n'utilise que le P\&L est un rétrécissement de l'historique de chaque trader vers
l'a priori du desk.}\label{lst:fm:hiring-and-compensation:shrink}

\begin{omfigure}
\begin{tikzpicture}
\begin{axis}[width=10.5cm,height=5.2cm,xlabel={\small ratio de Sharpe vrai},ylabel={\small rémunération sur cinq ans (millions de \$)},xmin=-0.4,xmax=2.0,ymin=0,
  tick label style={font=\footnotesize},grid=major,grid style={gray!20},table/col sep=comma]
\addplot[only marks,mark=*,mark size=1.4pt,omDef] table[x=sharpe,y=pay]{figdata/desk/10-hiring-and-compensation/scatter.csv};
\end{axis}
\end{tikzpicture}

\omcaption{Cinq ans de \omterm{def:fm:hiring-and-compensation:pool}{rémunération à la formule} (15~\% du P\&L annuel positif) face au ratio de Sharpe vrai du trader, pour 100 traders répartis sur cinq
desks simulés. La rémunération croît avec le talent, et des traders de même talent touchent entre presque rien et deux fois la moyenne. Données~:
\texttt{fm\_pay.one\_desk}.}\label{fig:fm:hiring-and-compensation:scatter}
\end{omfigure}

\begin{remark}[Pourquoi la rémunération à la formule persiste]\label{rem:fm:hiring-and-compensation:why}
Les \omterm{def:fm:hiring-and-compensation:pool}{rémunérations à la formule} sont fréquentes là où l'entreprise ne peut pas ou ne veut pas juger le talent, et là où le trader porte son propre drawdown (les plateformes
du chapitre~3)~: la formule est transparente et difficile à contester, et elle attire ceux qui croient à leur propre edge. Son coût~: elle rémunère la chance, et
elle rend la rémunération d'un trader convexe par rapport à son P\&L, ce qui récompense la prise de risque, sauf si des limites et le différé la plafonnent.
\end{remark}

\section{Différé, récupération et rétention}

Le différé remplit trois fonctions. Il expose la rémunération aux pertes qui apparaissent après l'année (par le biais du \omterm{def:fm:hiring-and-compensation:malus}{malus})~; il aligne la rémunération sur les résultats
de l'entreprise à plus long terme~; et il constitue un solde que le salarié perd en partant. Le plan du chapitre verse 60~\% de chaque attribution en numéraire et acquiert le reste en quatre
parts annuelles égales (\cref{fig:fm:hiring-and-compensation:vesting}). Un trader qui reçoit la même attribution chaque année détient, à chaque fin d'année, un
solde non acquis égal à une attribution annuelle complète~: $0.4\times(1+\tfrac34+\tfrac12+\tfrac14)=1$.

\begin{definition}[Rachat des différés]\label{def:fm:hiring-and-compensation:buyout}
Un \emph{rachat des différés}\index{rachat des différés} (deferral buyout) est le paiement qu'une entreprise d'embauche verse à un nouveau salarié pour remplacer la \omterm{def:fm:hiring-and-compensation:deferral}{rémunération différée} perdue en
quittant son ancien employeur, généralement sous forme d'instruments différés propres à la nouvelle entreprise et selon un calendrier semblable.
\end{definition}

\omcode{code/firm/bonuspool/firm_bonuspool.py}{59}{79}{Un plan de différé, les versements d'une attribution et le solde non acquis en fin d'année.}\label{lst:fm:hiring-and-compensation:defer}

\begin{omfigure}
\begin{tikzpicture}
\begin{axis}[width=10.5cm,height=5cm,ybar stacked,bar width=14pt,xlabel={\small année},ylabel={\small versé (en unités d'une attribution)},xmin=-0.6,xmax=7.6,
  ymin=0,ymax=1.1,xtick={0,1,2,3,4,5,6,7},tick label style={font=\footnotesize},grid=major,grid style={gray!20},table/col sep=comma,
  legend cell align=left,legend style={font=\scriptsize,at={(0.5,-0.3)},anchor=north,/tikz/every even column/.append style={column sep=6pt}},legend columns=5]
\addplot[fill=omDef!55,draw=omDef] table[x=year,y=c0]{figdata/desk/10-hiring-and-compensation/vesting.csv};
\addplot[fill=omProp!55,draw=omProp] table[x=year,y=c1]{figdata/desk/10-hiring-and-compensation/vesting.csv};
\addplot[fill=omThm!50,draw=omThm] table[x=year,y=c2]{figdata/desk/10-hiring-and-compensation/vesting.csv};
\addplot[fill=gray!40,draw=gray] table[x=year,y=c3]{figdata/desk/10-hiring-and-compensation/vesting.csv};
\addplot[fill=omDef!20,draw=omDef,dashed] table[x=year,y=c4]{figdata/desk/10-hiring-and-compensation/vesting.csv};
\legend{attribution 0,attribution 1,attribution 2,attribution 3,attribution 4}
\end{axis}
\end{tikzpicture}

\omcaption{Versements issus de cinq attributions annuelles égales dans un plan qui verse 60~\% en numéraire et acquiert 40~\% au prorata sur quatre ans. À partir de l'année 4, le
salarié reçoit une attribution complète par an, et dispose toujours d'une attribution complète non acquise. Données~: \texttt{firm.bonuspool.schedule}.}\label{fig:fm:hiring-and-compensation:vesting}
\end{omfigure}

Le solde est l'outil de rétention et le prix d'un changement d'employeur. Un concurrent qui embauche le trader doit racheter une attribution annuelle en
instruments différés qui lui sont propres, et le trader porte le risque de \omterm{def:fm:hiring-and-compensation:malus}{malus} de l'entreprise pendant les années où le solde s'acquiert. La \omterm{def:fm:hiring-and-compensation:malus}{récupération} va plus loin que le \omterm{def:fm:hiring-and-compensation:malus}{malus}, pour
reprendre ce qui a déjà été versé, et elle est plus difficile à faire valoir~: la règle américaine pour les dirigeants des sociétés cotées, les règles bancaires européennes et les plans de la plupart des entreprises privées
reposent surtout sur le \omterm{def:fm:hiring-and-compensation:malus}{malus}, appliqué à ce qui reste non acquis.

\section{Tutoriel~: vingt traders et trois règles de rémunération}\label{tut:fm:hiring-and-compensation:tut}

\begin{tutorial}
\textbf{Objectif.} Mesurer dans quelle mesure chaque règle de rémunération suit le talent, et la rétention que crée un plan de différé. \textbf{État final~:}
la \cref{fig:fm:hiring-and-compensation:years} et les corrélations du texte.
\begin{tutsteps}
\item \textbf{Le desk.} \texttt{fm\_pay.desk(rng)} tire les ratios de Sharpe vrais de vingt traders et cinq ans de P\&L.
\item \textbf{Les règles.} \texttt{fm\_pay.pay(pnl, rule)} applique \texttt{firm.bonuspool.formulaic}, \texttt{risk\_adjusted} et
\texttt{discretionary}.
\item \textbf{Les corrélations.} \texttt{fm\_pay.correlations()} calcule en moyenne, sur 2\,000 desks, la corrélation de la rémunération d'un an et de cinq ans avec le talent~;
\texttt{luck\_share\_formulaic()} calcule la part de chance d'une année de rémunération.
\item \textbf{Différé.} \texttt{firm.bonuspool.unvested([1.0]*6, DeferralPlan(0.6, 4), 5)} donne le solde non acquis en régime permanent, soit une attribution.
\end{tutsteps}
\end{tutorial}

\textbf{Pour aller plus loin.} Doubler la dispersion du talent et constater que le nombre d'années nécessaires baisse de quatre~; donner à la règle discrétionnaire une observation bruitée de
la qualité de décision de chaque trader, en plus du P\&L.

\section{Construction~: l'enveloppe de bonus}\label{bld:fm:hiring-and-compensation:bonuspool}

\begin{build}
\bfield{Rôle} L'arithmétique de rémunération du desk~: dimensionnement de l'enveloppe, règles d'allocation, différé et ses soldes, et part de la rémunération qui tient à la chance.
\bfield{Interface} \texttt{firm.bonuspool}~: \texttt{pool\_from\_profit}, \texttt{pool\_from\_revenue}~; \texttt{formulaic}, \texttt{risk\_adjusted},
\texttt{shrunk\_skill}, \texttt{discretionary}~; \texttt{DeferralPlan(cash\_share, years)}, \texttt{schedule}, \texttt{unvested}, \texttt{malus}~;
\texttt{luck\_share\_linear}, \texttt{years\_for\_correlation}.
\bfield{Règles} Une enveloppe n'est jamais négative~; une allocation discrétionnaire totalise exactement l'enveloppe~; les attributions différées s'acquièrent au prorata~; le \omterm{def:fm:hiring-and-compensation:malus}{malus} n'agit que sur le
solde non acquis.
\bfield{Tests d'acceptation} \texttt{code/firm/bonuspool/tests/}~: enveloppes et règles sur des chiffres établis à la main~; le poids du rétrécissement~; le calendrier et le
solde non acquis en régime permanent~; les formes closes de la \cref{prop:fm:hiring-and-compensation:years}.
\bfield{Pour aller plus loin} Rémunérer en parts de fonds dont la valeur suit le P\&L ultérieur du desk~; une attribution garantie la première année, et son coût~; l'enveloppe répartie
selon un crédit de Shapley (chapitre~9).
\end{build}

\begin{omsources}
\item F.~L.~Schmidt et J.~E.~Hunter, «~The validity and utility of selection methods in personnel psychology~», \emph{Psychological Bulletin}
124(2), 1998~; P.~R.~Sackett, C.~Zhang, C.~M.~Berry et F.~Lievens, «~Revisiting meta-analytic estimates of validity in personnel selection~»,
\emph{Journal of Applied Psychology} 107(11), 2022.
\item Directive 2013/36/EU, article 94, et Directive (EU) 2019/878~; PRA PS9/23 et FCA PS23/15 (2023)~; 17 CFR 240.10D-1.
\item European Banking Authority, communiqué de presse sur les hauts revenus de 2024, 16 avril 2026.
\end{omsources}

\section{Exercices}

\begin{exercise}[$\star$]\label{exo:fm:hiring-and-compensation:1}
Le P\&L annuel attendu d'un trader est de 8~millions de dollars, avec une volatilité de 10~millions de dollars. Quelles sont les probabilités d'une année à 20~millions de dollars ou plus et
d'une perte de 2~millions de dollars ou plus~?
\end{exercise}

\begin{exercise}[$\star$]\label{exo:fm:hiring-and-compensation:2}
Calculer les deux enveloppes de l'\cref{ex:fm:hiring-and-compensation:pools} pour une année de revenu de 120.
\end{exercise}

\begin{exercise}[$\star$]\label{exo:fm:hiring-and-compensation:3}
Un plan verse 60~\% de chaque attribution en numéraire et acquiert le reste au prorata sur quatre ans. Quel montant est versé au total la cinquième année à un trader qui a
reçu une attribution de 1 chaque année~?
\end{exercise}

\begin{exercise}[$\star\star$]\label{exo:fm:hiring-and-compensation:4}
Avec un écart de compétence $s=4$ et un bruit $n=10$, calculer la part de chance du P\&L d'une année et le nombre d'années nécessaires pour une corrélation de 0.8. Refaire le calcul pour
$s=8$.
\end{exercise}

\begin{exercise}[$\star\star$]\label{exo:fm:hiring-and-compensation:5}
Selon les règles européennes, la rémunération variable d'un \omterm{def:fm:hiring-and-compensation:mrt}{preneur de risques significatif} (material risk taker) atteint 150~\% de la rémunération fixe. Quelle approbation cela exige-t-il, et quel est le minimum qui doit être
différé et versé en instruments~?
\end{exercise}

\begin{exercise}[$\star\star$]\label{exo:fm:hiring-and-compensation:6}
Pourquoi une \omterm{def:fm:hiring-and-compensation:pool}{rémunération à la formule} récompense-t-elle la prise de risque, et quels sont les deux outils du chapitre qui la limitent~?
\end{exercise}

\begin{exercise}[$\star\star\star$]\label{exo:fm:hiring-and-compensation:7}
\emph{Programmation.} Relancer \texttt{fm\_pay.correlations(n\_desks=500, years=10)}. Quelle est la corrélation sur dix ans entre la \omterm{def:fm:hiring-and-compensation:pool}{rémunération à la formule} et la compétence~?
\end{exercise}

\begin{exercise}[$\star\star\star$]\label{exo:fm:hiring-and-compensation:8}
\emph{Trouver la faille.} «~Notre meilleur trader a gagné 25~millions de dollars l'an dernier, trois fois la moyenne du desk. Il faut lui offrir une somme garantie de 4~millions de dollars pour qu'il reste.~»
\end{exercise}

\section{Problème~: Payé pour la chance}

\begin{problem}[{Problème du week-end --- payé pour la chance}]\label{pb:fm:hiring-and-compensation:1}
Un responsable de desk doit choisir comment rémunérer vingt traders, et décider s'il faut égaler l'offre d'un concurrent faite à l'un d'entre eux.

\medskip\noindent\textbf{Partie~I --- Recrutement.}
\begin{enumerate}
\item Définir un \omterm{def:fm:hiring-and-compensation:structured}{entretien structuré} et dire ce que la recherche en conclut.
\item Concevoir un processus de recrutement en quatre étapes pour un trader et justifier son ordre.
\item Pourquoi un faible taux de base rend-il le premier tri important~?
\end{enumerate}

\medskip\noindent\textbf{Partie~II --- L'enveloppe.}
\begin{enumerate}[resume]
\item Définir une \omterm{def:fm:hiring-and-compensation:pool}{enveloppe de bonus} (bonus pool), une \omterm{def:fm:hiring-and-compensation:pool}{rémunération à la formule} et un \omterm{def:fm:hiring-and-compensation:pool}{bonus discrétionnaire}.
\item Dimensionner une enveloppe égale à 20~\% du profit après coûts et après \omterm{def:fm:limits-and-capital-allocation:raroc}{charge en capital}, pour un revenu de 100, des coûts de 60 et une charge de 10~; puis pour une année de revenu de 70.
\item Pourquoi une enveloppe calculée sur le profit après \omterm{def:fm:limits-and-capital-allocation:raroc}{charge en capital} diffère-t-elle d'une enveloppe calculée sur le revenu lors d'une mauvaise année~?
\end{enumerate}

\medskip\noindent\textbf{Partie~III --- Compétence et chance.}
\begin{enumerate}[resume]
\item Énoncer et démontrer la \cref{prop:fm:hiring-and-compensation:years}.
\item Donner la part de chance du P\&L d'une année et le nombre d'années nécessaires pour une corrélation de 0.8 sur le desk du chapitre.
\item Donner la part de chance de la \omterm{def:fm:hiring-and-compensation:pool}{rémunération à la formule} d'une année.
\item Donner la corrélation de chaque règle avec la compétence après un an et après cinq ans.
\item Pourquoi le pouvoir discrétionnaire fondé sur le seul P\&L ne bat-il pas la formule, et qu'est-ce qui le permettrait~?
\item Quelle est la corrélation sur dix ans entre la \omterm{def:fm:hiring-and-compensation:pool}{rémunération à la formule} et la compétence~?
\end{enumerate}

\medskip\noindent\textbf{Partie~IV --- Différé et offre.}
\begin{enumerate}[resume]
\item Définir la \omterm{def:fm:hiring-and-compensation:deferral}{rémunération différée}, un \omterm{def:fm:hiring-and-compensation:deferral}{calendrier d'acquisition}, le \omterm{def:fm:hiring-and-compensation:malus}{malus}, la \omterm{def:fm:hiring-and-compensation:malus}{récupération} (clawback) et le \omterm{def:fm:hiring-and-compensation:buyout}{rachat des différés} (buyout).
\item Donner le solde non acquis en régime permanent du plan du chapitre.
\item Que doit verser un concurrent pour embaucher un trader sur ce plan, et sous quelle forme~?
\item Énoncer les limites européennes concernant le ratio, le différé et les instruments, ainsi que la réforme britannique de 2023.
\item Pourquoi le \omterm{def:fm:hiring-and-compensation:malus}{malus} est-il plus facile à appliquer que la \omterm{def:fm:hiring-and-compensation:malus}{récupération}~?
\item Le responsable de desk doit-il égaler l'offre faite au trader qui a gagné 25~millions de dollars~?
\item Énoncer le \emph{résultat nommé}~: la part de chance d'une \omterm{def:fm:hiring-and-compensation:pool}{rémunération à la formule} sur un an, et le nombre d'années de P\&L au bout duquel la rémunération suit la compétence avec une
corrélation de 0.8.
\item En deux phrases, rédiger la politique de rémunération du desk.
\end{enumerate}
\end{problem}

\section{Questions d'entretien}

\begin{interviewq}[$\star$ \iqroles{trader}]\label{iq:fm:hiring-and-compensation:1}
Votre entreprise vous verse 15~\% de votre P\&L. Pourquoi pourrait-elle quand même en différer une partie~?
\end{interviewq}

\begin{interviewq}[$\star$ \iqroles{trader, chercheur}]\label{iq:fm:hiring-and-compensation:2}
Quelle est la différence entre le \omterm{def:fm:hiring-and-compensation:malus}{malus} et la \omterm{def:fm:hiring-and-compensation:malus}{récupération} (clawback)~?
\end{interviewq}

\begin{interviewq}[$\star\star$ \iqroles{chercheur}]\label{iq:fm:hiring-and-compensation:3}
Les ratios de Sharpe vrais des traders diffèrent avec un écart-type de 0.4, et le bruit de chaque P\&L annuel vaut 1 en unités de Sharpe. Combien d'années faut-il pour que le P\&L
classe les traders de façon fiable~?
\end{interviewq}

\begin{interviewq}[$\star\star$ \iqroles{trader}]\label{iq:fm:hiring-and-compensation:4}
Un concurrent vous propose un \omterm{def:fm:the-multi-manager-platform:payout}{taux de reversement} (payout) plus élevé, mais sans \omterm{def:fm:hiring-and-compensation:buyout}{rachat des différés} (buyout). Comment comparez-vous les offres~?
\end{interviewq}

\begin{interviewq}[$\star\star$ \iqroles{risque}]\label{iq:fm:hiring-and-compensation:5}
Pourquoi une \omterm{def:fm:hiring-and-compensation:pool}{rémunération à la formule} fondée sur le P\&L incite-t-elle à prendre des risques, et comment la contrôleriez-vous~?
\end{interviewq}

\begin{interviewq}[$\star\star\star$ \iqroles{chercheur, risque}]\label{iq:fm:hiring-and-compensation:6}
Comment concevriez-vous une règle de rémunération qui récompense la compétence plus vite que ne peut la révéler le seul P\&L~?
\end{interviewq}
